% ============================================================
%  CHAPTER 13: TEMPERATURE, THERMOMETERS AND THERMAL EXPANSION
%  The Physics Codex: A Complete University Foundation
%  Korede Samuel Omotosho (Falcon)
%  Aurikrex Learning Systems, 2026
%
%  File: chapters/part2/ch13_temperature_expansion.tex
% ============================================================

\chapter{Temperature, Thermometers and Thermal Expansion}
\label{ch:temperature_expansion}
\minitoc
\newpage

\chapterquote{Heat, like gravity, penetrates every
substance of the universe; its rays occupy all parts
of space.}{Joseph Fourier}

\begin{objectivebox}
By the end of this chapter, you should be able to:
\begin{enumerate}[noitemsep]
  \item Explain what temperature measures at the
  molecular level and distinguish it from heat
  \item Convert between Celsius, Fahrenheit and
  Kelvin scales
  \item Describe how different thermometers work
  and their thermometric properties
  \item State the desirable properties of a
  thermometric substance
  \item Calibrate a thermometer using fixed points
  and interpolation
  \item Calculate linear, area and volume expansion
  of solids and liquids
  \item Explain the bimetallic strip and its
  applications, including the thermostat
  \item Describe and explain the anomalous expansion
  of water and its consequences
  \item Distinguish between real and apparent
  expansion of a liquid
\end{enumerate}
\end{objectivebox}

% ============================================================
\section{Intuition: What is Temperature?}
% ============================================================

Place your hand on a metal table and then on a wooden
table that has been sitting in the same room. The metal
feels colder. Yet both are at exactly the same
temperature --- room temperature. Touch them with a
thermometer and both read the same value.

The difference is not temperature. The difference is
\textit{thermal conductivity} --- metal conducts heat
away from your hand faster than wood, so your hand
loses energy faster, and you interpret that faster
energy loss as coldness.

This reveals something important: your sense of ``hot''
and ``cold'' is not a reliable measure of temperature.
It measures heat flow, not temperature itself. For
that, we need instruments and precise definitions.

Temperature, at the molecular level, is a measure of
the \textbf{average kinetic energy} of the molecules
in a substance. A hot object has faster-moving molecules.
A cold object has slower-moving molecules. When two
objects at different temperatures are placed in contact,
energy flows from the faster (hotter) molecules to the
slower (colder) ones until the averages are equal ---
until \textbf{thermal equilibrium} is reached.

\begin{definitionbox}[Temperature, Heat and Thermal Equilibrium]
\textbf{Temperature} is a measure of the average
translational kinetic energy of the molecules in a
substance. It is a \textit{state property} of a body ---
it tells us how energetic the molecules are, and
determines the direction of heat flow.

\textbf{Heat} is energy in transit between two bodies
at different temperatures. It is not a property stored
in a body; it only exists during a transfer process.
A large block of ice contains more internal (thermal)
energy than a tiny hot spark, yet the spark has a far
higher temperature. Temperature and heat are
fundamentally different quantities.

\textbf{Thermal equilibrium} exists between two
objects when there is no net flow of heat between
them. They are at the same temperature.

\textbf{Zeroth Law of Thermodynamics:} If object A
is in thermal equilibrium with object C, and object
B is also in thermal equilibrium with C, then A and
B are in thermal equilibrium with each other.
(This law justifies the use of thermometers: a
thermometer in equilibrium with your body reads your
body temperature.)
\end{definitionbox}

% ============================================================
\section{Temperature Scales}
% ============================================================

Three temperature scales are in common use in science
and everyday life.

\begin{definitionbox}[The Three Temperature Scales]
\textbf{Celsius scale (°C):}
\begin{itemize}[noitemsep]
  \item Lower fixed point (ice point): $0\,°\text{C}$
  (pure melting ice and water in equilibrium at 1 atm)
  \item Upper fixed point (steam point): $100\,°\text{C}$
  (pure water and steam in equilibrium at 1 atm)
\end{itemize}

\textbf{Kelvin scale (K)} --- the SI base unit of
thermodynamic temperature:
\begin{itemize}[noitemsep]
  \item Absolute zero: $0\ \text{K}$ (theoretical
  minimum --- all molecular motion ceases; unattainable
  in practice per the Third Law of Thermodynamics)
  \item Triple point of water: $273.16\ \text{K}$
  \item $0\,°\text{C} = 273.15\ \text{K}
  \approx 273\ \text{K}$
\end{itemize}
\[
  T(\text{K}) = T(°\text{C}) + 273.15
\]
\begin{notebox}
The Kelvin is a unit, not a degree. Write $300\ \text{K}$,
\textbf{not} $300\,°\text{K}$. The degree symbol is
never used with Kelvin. Note also that a temperature
\emph{difference} of $1\,°\text{C}$ equals a difference
of $1\ \text{K}$; only the zero points differ.
\end{notebox}

\textbf{Fahrenheit scale (°F):}
\[
  T(°\text{F}) = \frac{9}{5}T(°\text{C}) + 32
\]
\[
  T(°\text{C}) = \frac{5}{9}\left[T(°\text{F}) - 32\right]
\]
For direct conversion from Fahrenheit to Kelvin,
combine both relationships:
\[
  T(\text{K}) = \frac{5}{9}\left[T(°\text{F})
  + 459.67\right] \approx \frac{5}{9}
  \left[T(°\text{F}) + 460\right]
\]
\end{definitionbox}

\begin{diagbox}[Comparing Temperature Scales]
\begin{center}
\begin{tikzpicture}[scale=1.0, font=\small]

  % Three vertical axes
  % Kelvin axis
  \draw[very thick, codexblue] (1,0) -- (1,9);
  \node[codexblue, font=\bfseries] at (1,9.5) {Kelvin (K)};

  % Celsius axis
  \draw[very thick, codexred] (5,0) -- (5,9);
  \node[codexred, font=\bfseries] at (5,9.5) {Celsius (°C)};

  % Fahrenheit axis
  \draw[very thick, codexgold] (9,0) -- (9,9);
  \node[codexgold, font=\bfseries] at (9,9.5) {Fahrenheit (°F)};

  % Key temperatures with tick marks and connectors
  % --- Absolute zero ---
  \draw[thick, codexblue] (0.7,0.5) -- (1.3,0.5);
  \node[left, codexblue] at (0.7,0.5) {0 K};
  \draw[thick, codexred] (4.7,0.5) -- (5.3,0.5);
  \node[right, codexred] at (5.3,0.5) {$-273.15°$C};
  \draw[thick, codexgold] (8.7,0.5) -- (9.3,0.5);
  \node[right, codexgold] at (9.3,0.5) {$-459.67°$F};
  \draw[dashed, gray] (1.3,0.5) -- (8.7,0.5);

  % --- Ice point ---
  \draw[thick, codexblue] (0.7,3.5) -- (1.3,3.5);
  \node[left, codexblue] at (0.7,3.5) {273.15 K};
  \draw[thick, codexred] (4.7,3.5) -- (5.3,3.5);
  \node[right, codexred] at (5.3,3.5) {0°C};
  \draw[thick, codexgold] (8.7,3.5) -- (9.3,3.5);
  \node[right, codexgold] at (9.3,3.5) {32°F};
  \draw[dashed, gray] (1.3,3.5) -- (8.7,3.5);
  \node[above, gray, font=\tiny] at (5,3.6)
    {Ice point};

  % --- Body temperature ---
  \draw[thick, codexblue] (0.7,5.2) -- (1.3,5.2);
  \node[left, codexblue] at (0.7,5.2) {310.15 K};
  \draw[thick, codexred] (4.7,5.2) -- (5.3,5.2);
  \node[right, codexred] at (5.3,5.2) {37°C};
  \draw[thick, codexgold] (8.7,5.2) -- (9.3,5.2);
  \node[right, codexgold] at (9.3,5.2) {98.6°F};
  \draw[dashed, gray] (1.3,5.2) -- (8.7,5.2);
  \node[above, gray, font=\tiny] at (5,5.3)
    {Body temp.};

  % --- Steam point ---
  \draw[thick, codexblue] (0.7,7.0) -- (1.3,7.0);
  \node[left, codexblue] at (0.7,7.0) {373.15 K};
  \draw[thick, codexred] (4.7,7.0) -- (5.3,7.0);
  \node[right, codexred] at (5.3,7.0) {100°C};
  \draw[thick, codexgold] (8.7,7.0) -- (9.3,7.0);
  \node[right, codexgold] at (9.3,7.0) {212°F};
  \draw[dashed, gray] (1.3,7.0) -- (8.7,7.0);
  \node[above, gray, font=\tiny] at (5,7.1)
    {Steam point};

  % Scale markings (K has larger intervals)
  \foreach \y in {0.5,1.0,...,8.5} {
    \draw[gray!30] (0.8,\y) -- (1.2,\y);
  }

\end{tikzpicture}
\end{center}
\end{diagbox}

\begin{examplebox}[13.1]
\stepcounter{workedexample}
Convert the following temperatures:\\
(a) $37\,°\text{C}$ (body temperature) to Kelvin and
Fahrenheit.\\
(b) $-40\,°\text{F}$ to Celsius (this is the one
temperature where the two scales coincide).\\
(c) The surface of the Sun is approximately
$5778\ \text{K}$. Express in Celsius.
\end{examplebox}

\begin{solutionbox}
\textbf{(a)}
$T(\text{K}) = 37 + 273 = 310\ \text{K}$\\
$T(°\text{F}) = \frac{9}{5}(37) + 32 = 66.6 + 32
= 98.6\,°\text{F}$

\textbf{(b)}
$T(°\text{C}) = \frac{5}{9}(-40 - 32)
= \frac{5}{9}(-72) = -40\,°\text{C}$\\
Indeed, $-40\,°\text{F} = -40\,°\text{C}$ --- the unique
crossing point of the two scales.

\textbf{(c)}
$T(°\text{C}) = 5778 - 273 = 5505\,°\text{C}$\\
The surface of the Sun is about $5500\,°\text{C}$.
\end{solutionbox}

% ============================================================
\section{Thermometers and Thermometric Properties}
% ============================================================

\begin{definitionbox}[Thermometric Property]
A \textbf{thermometric property} is any physical
property that changes measurably, continuously and
reproducibly with temperature and can be used as the
basis for a thermometer.

\textbf{Desirable properties of a thermometric
substance:}
\begin{itemize}[noitemsep]
  \item \textbf{Sensitivity:} the property should
  change sufficiently for small changes in temperature
  to be detectable.
  \item \textbf{Linearity:} ideally, the property
  varies linearly with temperature, so that
  interpolation between fixed points is accurate.
  \item \textbf{Reproducibility:} the same temperature
  always produces the same value of the property.
  \item \textbf{Wide range:} the thermometric substance
  should remain in a usable state over as wide a
  temperature range as possible.
  \item \textbf{Response time:} the thermometer should
  reach thermal equilibrium with the body being
  measured quickly.
  \item \textbf{Safety and practicality:} the substance
  should be non-corrosive, non-toxic if possible, and
  easy to work with.
\end{itemize}
\end{definitionbox}

\begin{table}[H]
\centering
\caption{Types of Thermometer and Their Thermometric Properties}
\label{tab:thermometers}
\begin{tabular}{llll}
\toprule
\textbf{Thermometer} & \textbf{Property} & \textbf{Range} &
\textbf{Application}\\
\midrule
Liquid-in-glass & Volume of liquid &
  $-10$ to $110\,°$C & Everyday use\\
Resistance (platinum) & Electrical resistance &
  $-200$ to $1200\,°$C & Precision measurement\\
Thermocouple & EMF (voltage) &
  $-250$ to $1400\,°$C & Industry, engines\\
Gas thermometer & Pressure/volume of gas &
  $-270$ to $1500\,°$C & Standard reference\\
Pyrometer & Radiation intensity &
  $>600\,°$C & Furnaces, stars\\
Clinical & Mercury/galinstan volume &
  $35$ to $42\,°$C & Medical\\
\bottomrule
\end{tabular}
\end{table}

\subsection{The Liquid-in-Glass Thermometer}

\begin{diagbox}[Liquid-in-Glass Thermometer: Cross-Section]
\begin{center}
\begin{tikzpicture}[scale=1.0, font=\small]

  % Main tube
  \fill[gray!20] (0,0.3) rectangle (8.5,0.7);
  \draw[thick] (0,0.3) rectangle (8.5,0.7);

  % Capillary tube (inner)
  \fill[white] (0.5,0.38) rectangle (8.0,0.62);
  \draw[thick] (0.5,0.38) -- (8.0,0.38);
  \draw[thick] (0.5,0.62) -- (8.0,0.62);

  % Liquid column (mercury or alcohol)
  \fill[codexred!60] (0.5,0.38) rectangle (4.5,0.62);

  % Bulb at left end
  \fill[codexred!60] (0,0.0) .. controls (-0.3,0.0) and
    (-0.3,1.0) .. (0,1.0) -- (0.5,1.0) --
    (0.5,0.0) -- cycle;
  \draw[thick] (0.5,0.0) -- (0,0.0) .. controls
    (-0.3,0.0) and (-0.3,1.0) .. (0,1.0) -- (0.5,1.0);

  % Sealed end at right
  \draw[thick] (8.5,0.3) arc(-90:90:0.2);

  % Scale markings
  \foreach \x in {1,2,...,7} {
    \draw[thick] (\x+0.5, 0.62) -- (\x+0.5, 0.75);
  }
  \node[above, font=\tiny] at (1.5,0.75) {10};
  \node[above, font=\tiny] at (3.5,0.75) {30};
  \node[above, font=\tiny] at (5.5,0.75) {50};
  \node[above, font=\tiny] at (7.5,0.75) {70};

  % Labels
  \node[below] at (0.2,0.0) {\small Bulb};
  \node[below] at (4.0,-0.3) {\small Liquid column};
  \node[below] at (7.5,0.3) {\small Capillary tube};

  % Constriction (clinical thermometer feature)
  \fill[white] (2.4,0.38) rectangle (2.6,0.62);
  \fill[gray!30] (2.45,0.38) rectangle (2.55,0.62);
  \draw[thick] (2.45,0.38) -- (2.45,0.62);
  \draw[thick] (2.55,0.38) -- (2.55,0.62);
  \node[below, font=\tiny, align=center] at (2.5,0.3)
    {\small constriction\\(clinical only)};

  % Mercury label
  \node[codexred, font=\small] at (2.0,1.3)
    {Mercury (or coloured alcohol)};
  \draw[->, codexred] (2.0,1.15) -- (2.0,0.65);

\end{tikzpicture}
\end{center}
\end{diagbox}

\begin{notebox}[Mercury vs.\ Alcohol]
\textbf{Mercury} is preferred for general-purpose
thermometers because it is a good thermal conductor,
does not wet (stick to) glass, and remains liquid
from $-38.8\,°\text{C}$ to $356.7\,°\text{C}$.
Its expansion is nearly linear with temperature.
However, mercury is toxic, and its use in clinical
thermometers is being phased out in many countries
in favour of \textit{galinstan} (a non-toxic gallium
alloy) or digital sensors.

\textbf{Alcohol} (coloured red or blue with dye)
stays liquid down to about $-115\,°\text{C}$,
making it essential for low-temperature measurements
--- far below mercury's freezing point. Its main
drawbacks are a low boiling point ($\approx 78\,°\text{C}$
for ethanol), which limits its upper range, and the
fact that alcohol wets glass and tends to separate
in the tube if cooled too rapidly.
\end{notebox}

\subsection{Calibration and Scale Interpolation}

\begin{processbox}[Calibrating a Thermometer]
To calibrate any thermometer:
\begin{enumerate}[noitemsep]
  \item Immerse in melting pure ice at 1 atm.
  Record the thermometric property value $X_0$.
  Mark this as $0\,°\text{C}$.
  \item Immerse in steam above boiling water at 1 atm.
  Record the property value $X_{100}$.
  Mark this as $100\,°\text{C}$.
  \item Interpolate: for any property value $X_\theta$:
\end{enumerate}
\[
  \theta = \frac{X_\theta - X_0}{X_{100} - X_0}
  \times 100\,°\text{C}
\]
This formula assumes the thermometric property varies
\textit{linearly} with temperature between the two
fixed points. Different thermometric properties do
not, in general, vary in exactly the same way, so
two thermometers calibrated by this method may give
slightly different readings between the fixed points.
The \textit{ideal gas thermometer} is the primary
standard because an ideal gas expands perfectly
linearly (by definition of the thermodynamic
temperature scale).
\end{processbox}

\begin{diagbox}[Thermometer Calibration]
\begin{center}
\begin{tikzpicture}[scale=1.0, font=\small]

  % Axes
  \draw[->, thick] (0,0) -- (9,0)
    node[right]{Thermometric property $X$};
  \draw[->, thick] (0,0) -- (0,7)
    node[above]{Temperature $\theta$ (°C)};

  % Key points
  \fill[codexblue] (1.5,0) circle (4pt);
  \fill[codexblue] (1.5,1.0) circle (2pt);
  \draw[dashed, codexblue] (1.5,0) -- (1.5,1.0)
    -- (0,1.0);
  \node[below, codexblue] at (1.5,0) {$X_0$};
  \node[left, codexblue] at (0,1.0) {$0\,°$C};

  \fill[codexred] (7.0,0) circle (4pt);
  \fill[codexred] (7.0,6.0) circle (2pt);
  \draw[dashed, codexred] (7.0,0) -- (7.0,6.0)
    -- (0,6.0);
  \node[below, codexred] at (7.0,0) {$X_{100}$};
  \node[left, codexred] at (0,6.0) {$100\,°$C};

  % Unknown point
  \fill[codexgold] (4.3,0) circle (4pt);
  \fill[codexgold] (4.3,3.5) circle (3pt);
  \draw[dashed, codexgold] (4.3,0) -- (4.3,3.5)
    -- (0,3.5);
  \node[below, codexgold] at (4.3,0) {$X_\theta$};
  \node[left, codexgold] at (0,3.5) {$\theta$};

  % Linear relationship line
  \draw[very thick, codexblue]
    (1.5,1.0) -- (7.0,6.0);

  % Formula annotation
  \node[align=center, font=\small, codexblue]
    at (10.0,3.0)
    {$\theta = \dfrac{X_\theta - X_0}
    {X_{100} - X_0} \times 100\,°\text{C}$};


\end{tikzpicture}
\end{center}
\end{diagbox}

\begin{examplebox}[13.2]
\stepcounter{workedexample}
A platinum resistance thermometer has resistance
$25.0\ \Omega$ at the ice point and $35.4\ \Omega$
at the steam point. At an unknown temperature, the
resistance is $29.6\ \Omega$.\\
(a) Calculate the unknown temperature on the
platinum scale.\\
(b) If a mercury thermometer simultaneously reads
$43.5\,°\text{C}$, explain why there is a discrepancy
and which reading is more reliable.
\end{examplebox}

\begin{solutionbox}
\textbf{(a)}
\[
  \theta = \frac{X_\theta - X_0}{X_{100} - X_0}
  \times 100
  = \frac{29.6 - 25.0}{35.4 - 25.0} \times 100
  = \frac{4.6}{10.4} \times 100
  = 44.2\,°\text{C}
\]

\textbf{(b)} Different thermometers use different
thermometric properties. The interpolation formula
assumes a linear relationship between the property
and temperature, but this is only truly satisfied by
the ideal gas thermometer (which defines the
thermodynamic scale). Platinum resistance and mercury
volume both change with temperature, but not in
exactly the same way across all temperatures. The
discrepancy ($44.2$ vs $43.5\,°$C) reflects this
non-linearity in one or both properties. The platinum
resistance thermometer is generally more reliable,
more stable, and more reproducible than the
liquid-in-glass thermometer, and it covers a far
wider temperature range.
\end{solutionbox}

% ============================================================
\section{Thermal Expansion of Solids}
% ============================================================

When a solid is heated, the molecules vibrate more
vigorously about their equilibrium positions. Due to
the asymmetry of intermolecular potential energy
curves, the \textit{average} separation between
molecules increases. The solid expands.

\begin{diagbox}[Molecular Origin of Thermal Expansion]
\begin{center}
\begin{tikzpicture}[scale=1.0, font=\small]
  % Potential energy vs separation curve
  \begin{axis}[
    width=10cm, height=6cm,
    xlabel={Molecular separation $r$},
    ylabel={Potential energy $E_p$},
    xmin=0.5, xmax=3.5,
    ymin=-1.3, ymax=1.0,
    ytick={-1.0,0},
    yticklabels={min $E_p$,0},
    xtick={1.2},
    xticklabels={$r_0$},
    grid=major, grid style={gray!15},
    clip=true,
    restrict y to domain=-1.3:1.0,
    unbounded coords=jump
  ]
    % Potential energy curve (Lennard-Jones like)
    \addplot[codexblue, very thick, domain=0.7:3.5,
      samples=100]
      {4*(1/x^12 - 1/x^6) + 0.2};
    % Equilibrium position
    \addplot[dashed, codexgold, domain=0.7:3.5]{0.2-1.0};
    \fill[codexgold] (axis cs:1.12,-0.8) circle (3pt);
    \draw[dashed, codexgold, thick]
      (axis cs:1.12,-2) -- (axis cs:1.12,-0.8);
    % Low temperature energy level
    \addplot[codexgreen, thick, domain=1.063:1.218]{-0.65};
    \fill[codexgreen] (axis cs:1.063,-0.65) circle (2pt);
    \fill[codexgreen] (axis cs:1.218,-0.65) circle (2pt);
    \node[codexgreen, right, align=left, font=\tiny]
      at (axis cs:1.24,-0.65) {Low $T$\\avg $r_1$};
    % High temperature energy level
    \addplot[codexred, thick, domain=1.027:1.377]{-0.3};
    \fill[codexred] (axis cs:1.027,-0.3) circle (2pt);
    \fill[codexred] (axis cs:1.377,-0.3) circle (2pt);
    \node[codexred, right, align=left, font=\tiny]
      at (axis cs:1.40,-0.3) {High $T$\\avg $r_2 > r_1$};
    % Arrows showing average positions (own well-separated rows below the frame)
    \draw[<->, codexgreen, dashed]
      (axis cs:1.063,-1.15) -- (axis cs:1.218,-1.15);
    \node[codexgreen, below, font=\tiny]
      at (axis cs:1.1405,-1.15) {$\langle r \rangle_1$};
    \draw[<->, codexred, dashed]
      (axis cs:1.027,-1.25) -- (axis cs:1.377,-1.25);
    \node[codexred, below, font=\tiny]
      at (axis cs:1.202,-1.25) {$\langle r \rangle_2$};
    % Note: asymmetric curve causes net expansion
    \node[align=center, font=\tiny, gray]
      at (axis cs:2.8,0.6)
      {Asymmetric curve:\\left side steep\\right side gentle\\
      $\Rightarrow$ avg separation\\increases with energy};
  \end{axis}
\end{tikzpicture}
\end{center}
\small The asymmetry of the potential energy curve
means that as energy increases (higher temperature),
the average separation increases --- the solid expands.
If the curve were symmetric, the average position
would not shift and there would be no thermal expansion.
\end{diagbox}

\subsection{Linear Expansion}

\begin{lawbox}[Linear Expansion]
When a solid rod of original length $L_0$ is heated
through a temperature change $\Delta\theta$, its
length increases by:
\[
  \Delta L = \alpha L_0 \Delta\theta
\]
\[
  L = L_0(1 + \alpha\Delta\theta)
\]
where $\alpha$ is the \textbf{coefficient of linear
expansion} (also written $\alpha_L$). Unit:
$\text{K}^{-1}$ or $°\text{C}^{-1}$.

\begin{notebox}
Strictly, $\alpha$ is itself temperature-dependent.
The formula $\Delta L = \alpha L_0 \Delta\theta$
treats $\alpha$ as constant, which is an excellent
approximation for small temperature ranges
($\Delta\theta \lesssim 100\,\text{K}$) and the
temperature ranges encountered in most engineering
problems. For very large temperature changes, a
temperature-averaged value of $\alpha$ should be used.
\end{notebox}
\end{lawbox}

\begin{table}[H]
\centering
\caption{Coefficients of Expansion of Common Materials}
\label{tab:expansion}
\begin{tabular}{lll}
\toprule
\textbf{Material} & \textbf{Linear $\alpha$
(×10$^{-6}$ K$^{-1}$)} & \textbf{Volume $\gamma$
(×10$^{-6}$ K$^{-1}$)}\\
\midrule
Invar (steel alloy) & 1.2 & 3.6\\
Glass (Pyrex) & 3.2 & 9.6\\
Concrete & 12 & 36\\
Steel & 12 & 36\\
Copper & 17 & 51\\
Brass & 19 & 57\\
Aluminium & 23 & 69\\
Lead & 29 & 87\\
Mercury (liquid) & --- & 182\\
Water at 20°C & --- & 207\\
\bottomrule
\end{tabular}
\end{table}

\subsection{Area and Volume Expansion}

\begin{lawbox}[Area and Volume Expansion]
\textbf{Area expansion:}
\[
  \Delta A = \beta A_0 \Delta\theta, \quad
  A = A_0(1 + \beta\Delta\theta), \quad
  \beta = 2\alpha
\]

\textbf{Volume expansion:}
\[
  \Delta V = \gamma V_0 \Delta\theta, \quad
  V = V_0(1 + \gamma\Delta\theta), \quad
  \gamma = 3\alpha
\]
where $\beta$ is the coefficient of area expansion
and $\gamma$ is the coefficient of volume expansion.

\textbf{Derivation of $\beta = 2\alpha$:}
For a square of side $L$, area $A_0 = L^2$:
\[
  A = L^2(1+\alpha\Delta\theta)^2
  \approx L^2(1+2\alpha\Delta\theta)
  = A_0(1+2\alpha\Delta\theta)
\]
(ignoring $\alpha^2(\Delta\theta)^2$ since $\alpha$
is of order $10^{-5}\ \text{K}^{-1}$, so this term
is negligibly small)

\textbf{Key insight --- holes expand too:}
A hole in a solid expands just as if it were filled
with the same solid material. This is because
thermal expansion is uniform --- every linear
dimension scales by the same factor
$(1 + \alpha\Delta\theta)$.
\end{lawbox}

\begin{diagbox}[Linear, Area and Volume Expansion Illustrated]
\begin{center}
\begin{tikzpicture}[scale=0.9, font=\small]
  % === LINEAR ===
  \begin{scope}[shift={(0,5)}]
    \node[font=\bfseries] at (2.5,1.5) {Linear ($\alpha$)};
    % Before
    \fill[codexblue!30] (0,0.3) rectangle (2.5,0.7);
    \draw[thick] (0,0.3) rectangle (2.5,0.7);
    \draw[<->, gray] (0,-0.1) -- (2.5,-0.1);
    \node[below, gray] at (1.25,-0.15) {$L_0$};
    % Arrow
    \draw[->, very thick, codexred] (3.0,0.5) -- (3.7,0.5)
      node[above, font=\tiny]{heat};
    % After
    \fill[codexred!30] (4.0,0.3) rectangle (7.0,0.7);
    \draw[thick] (4.0,0.3) rectangle (7.0,0.7);
    \draw[<->, gray] (4.0,-0.1) -- (7.0,-0.1);
    \node[below, gray, font=\tiny] at (5.5,-0.15) {$L_0(1+\alpha\Delta\theta)$};
    % Extension marker, with thin connector guides tying the
    % ΔL bracket visually back to the highlighted segment it labels
    \fill[codexred!60] (6.4,0.3) rectangle (7.0,0.7);
    \draw[codexred!50, thin] (6.4,0.3) -- (6.4,-0.55);
    \draw[codexred!50, thin] (7.0,0.3) -- (7.0,-0.55);
    \draw[<->, codexred] (6.4,-0.55) -- (7.0,-0.55);
    \node[below, codexred, font=\tiny] at (6.7,-0.65)
      {$\Delta L$};
  \end{scope}
  % === AREA ===
  \begin{scope}[shift={(0,1.5)}]
    \node[font=\bfseries] at (2.0,2.3) {Area ($\beta = 2\alpha$)};
    % Before
    \fill[codexblue!30] (0,0) rectangle (1.5,1.5);
    \draw[thick] (0,0) rectangle (1.5,1.5);
    \node[font=\tiny] at (0.75,0.75) {$A_0$};
    % Arrow
    \draw[->, very thick, codexred] (2.0,0.75) -- (2.7,0.75)
      node[above, font=\tiny]{heat};
    % After
    \fill[codexblue!20] (3.0,0) rectangle (4.8,1.8);
    \fill[codexred!30] (3.0,0) rectangle (4.8,1.8);
    \draw[thick] (3.0,0) rectangle (4.8,1.8);
    \node[font=\tiny] at (3.9,2.0) {$A_0(1+\beta\Delta\theta)$};
    % Extension markers
    \draw[<->, codexred] (3.0,-0.3) -- (4.8,-0.3);
    \node[below, codexred, font=\tiny] at (3.9,-0.4)
      {expands in both dimensions};
  \end{scope}
  % === VOLUME ===
  \begin{scope}[shift={(6,1.0)}]
    \node[font=\bfseries] at (2.2,2.8) {Volume ($\gamma = 3\alpha$)};
    % Cube before (3D)
    \fill[codexblue!30] (0,0) -- (1.5,0) -- (1.5,1.5) -- (0,1.5) -- cycle;
    \fill[codexblue!20] (0,1.5) -- (0.5,2.0) -- (2.0,2.0) -- (1.5,1.5) -- cycle;
    \fill[codexblue!25] (1.5,0) -- (2.0,0.5) -- (2.0,2.0) -- (1.5,1.5) -- cycle;
    \draw[thick] (0,0)--(1.5,0)--(1.5,1.5)--(0,1.5)--(0,0);
    \draw[thick] (0,1.5)--(0.5,2.0)--(2.0,2.0)--(1.5,1.5);
    \draw[thick] (1.5,0)--(2.0,0.5)--(2.0,2.0);
    \node[font=\tiny] at (0.75,0.75) {$V_0$};
    % Arrow
    \draw[->, very thick, codexred]
      (2.3,1.0) -- (2.75,1.0)
      node[above, font=\tiny]{heat};
    % Cube after (larger, 3D)
    \fill[codexred!20] (3.2,0) -- (5.2,0) -- (5.2,2.0) -- (3.2,2.0) -- cycle;
    \fill[codexred!10] (3.2,2.0) -- (3.8,2.6) -- (5.8,2.6) -- (5.2,2.0) -- cycle;
    \fill[codexred!15] (5.2,0) -- (5.8,0.6) -- (5.8,2.6) -- (5.2,2.0) -- cycle;
    \draw[thick] (3.2,0)--(5.2,0)--(5.2,2.0)--(3.2,2.0)--(3.2,0);
    \draw[thick] (3.2,2.0)--(3.8,2.6)--(5.8,2.6)--(5.2,2.0);
    \draw[thick] (5.2,0)--(5.8,0.6)--(5.8,2.6);
    \node[font=\tiny] at (4.2,1.0) {$V_0(1+\gamma\Delta\theta)$};
  \end{scope}
\end{tikzpicture}
\end{center}
\end{diagbox}

\begin{examplebox}[13.3]
\stepcounter{workedexample}
A steel railway track has a length of $25.0\ \text{m}$
at $20\,°\text{C}$.
($\alpha_{\text{steel}} = 12\times10^{-6}\ \text{K}^{-1}$)\\
(a) Calculate its length at $55\,°\text{C}$.\\
(b) Calculate the expansion.\\
(c) What gap must be left between tracks so they
do not buckle when the temperature rises from
$10\,°\text{C}$ (winter) to $55\,°\text{C}$ (hot
Nigerian afternoon)?
\end{examplebox}

\begin{solutionbox}
\textbf{(a)} Length at $55\,°\text{C}$:
\[
  \Delta\theta = 55 - 20 = 35\ \text{K}
\]
\[
  L = L_0(1 + \alpha\Delta\theta)
  = 25.0\left(1 + 12\times10^{-6} \times 35\right)
\]
\[
 = 25.0(1 + 4.2\times10^{-4})
  = 25.0 \times 1.00042
  = 25.0105\ \text{m}
\]

\textbf{(b)} Expansion:
$\Delta L = 25.0105 - 25.0 = 0.0105\ \text{m}
= 10.5\ \text{mm}$

\textbf{(c)} Temperature range: $\Delta\theta = 55 - 10
= 45\ \text{K}$
\[
  \Delta L = \alpha L_0 \Delta\theta
  = 12\times10^{-6} \times 25.0 \times 45
  = 0.0135\ \text{m} = 13.5\ \text{mm}
\]
The expansion gap must be at least $13.5\ \text{mm}$
per $25\ \text{m}$ track section. Modern rail systems
use \textit{continuously welded rail} (CWR) where the
track is held under tension, but the calculation
still underpins the design specification. Similar
principles apply to expansion joints in long bridges,
concrete highways, and building structures.
\end{solutionbox}

% ============================================================
\section{The Bimetallic Strip}
% ============================================================

\begin{diagbox}[The Bimetallic Strip: Principle and Applications]
\begin{center}
\begin{tikzpicture}[scale=1.0, font=\small]

  % === AT ROOM TEMPERATURE ===
  \begin{scope}[shift={(0,4.5)}]
    \node[font=\bfseries] at (4,1.8) {At room temperature: straight};
    % Brass strip (top, higher $\alpha$)
    \fill[codexgold!60] (0,0.5) rectangle (7,1.0);
    \draw[thick] (0,0.5) rectangle (7,1.0);
    \node[codexgold, font=\tiny] at (3.5,0.75) {Brass ($\alpha$ larger)};
    % Steel strip (bottom, lower $\alpha$)
    \fill[gray!40] (0,0) rectangle (7,0.5);
    \draw[thick] (0,0) rectangle (7,0.5);
    \node[gray, font=\tiny] at (3.5,0.25) {Steel ($\alpha$ smaller)};
    % Fixed end
    \fill[gray!70] (-0.3,-0.2) rectangle (0,1.2);
    \draw[thick] (-0.3,-0.2) rectangle (0,1.2);
    \node[below, font=\tiny] at (-0.15,-0.3) {fixed};
  \end{scope}

  % === WHEN HEATED ===
  \begin{scope}[shift={(0,1.5)}]
    \node[font=\bfseries] at (4.5,2.3) {When heated: curves downward (brass expands more)};
    % Fixed end
    \fill[gray!70] (-0.3,-0.2) rectangle (0,1.2);
    \draw[thick] (-0.3,-0.2) rectangle (0,1.2);
    % Curved brass (top, longer, on outside of curve)
    \draw[codexgold!80, very thick, fill=codexgold!50]
      (0,0.5) .. controls (2,0.9) and (5,1.4) ..
      (7,2.0) -- (7,1.5) ..
      controls (5,0.9) and (2,0.4) ..
      (0,0.0) -- cycle;
    \node[codexgold, font=\tiny] at (3.5,1.0)
      {Brass (outer, longer)};
    % Curved steel (bottom, shorter, inside of curve)
    \draw[gray!60, very thick, fill=gray!40]
      (0,0.0) .. controls (2,0.4) and (5,0.9) ..
      (7,1.5) -- (7,1.1) ..
      controls (5,0.5) and (2,0.1) ..
      (0,-0.4) -- cycle;
    \node[gray, font=\tiny] at (3.5,0.5)
      {Steel (inner, shorter)};
    % Curvature arrow
    \draw[->, very thick, codexred]
      (7.2,1.5) -- (7.2,0.5)
      node[right, align=center]{bends\\down};
  \end{scope}

  % === APPLICATIONS ===
  \begin{scope}[shift={(0,-1.5)}]
    \node[font=\bfseries] at (4.5,2.3) {Applications:};
    \node[font=\small, align=left] at (4.0,0.9)
      {$\bullet$ Thermostats (heating/cooling control)\\
       $\bullet$ Fire alarms (circuit breaks when heated)\\
       $\bullet$ Electrical circuit breakers\\
       $\bullet$ Oven temperature indicators\\
       $\bullet$ Bimetallic thermometers};
  \end{scope}

\end{tikzpicture}
\end{center}
\end{diagbox}

\begin{realworldbox}[The Bimetallic Thermostat]
In a simple thermostat, a bimetallic strip acts as a
temperature-sensitive switch. At room temperature
(below the set point), the strip is straight and
presses against an electrical contact, completing
the circuit that powers the heater. As the room
warms above the set point, the strip bends away from
the contact, breaking the circuit and switching the
heater off. When the room cools again, the strip
straightens and re-makes the contact. This
on--off cycling maintains the room near the desired
temperature. An adjusting screw moves the contact
point closer or further, setting the trigger
temperature. More modern thermostats use electronic
sensors and relays, but bimetallic strips remain
common in simple, inexpensive applications such as
electric kettles, toasters, and car engine cooling
fans.
\end{realworldbox}

% ============================================================
\section{Expansion of Liquids}
% ============================================================

Liquids only expand in volume (they have no fixed shape).
The apparent expansion of a liquid in a container is
less than its real expansion because the container also
expands.

\begin{lawbox}[Real and Apparent Expansion of a Liquid]
\[
  \gamma_{\text{real}} = \gamma_{\text{apparent}}
  + \gamma_{\text{container}}
\]
where:
\begin{itemize}[noitemsep]
  \item $\gamma_{\text{real}}$ = true coefficient of
  volume expansion of the liquid
  \item $\gamma_{\text{apparent}}$ = observed coefficient
  (expansion of liquid relative to the container)
  \item $\gamma_{\text{container}}$ = volume expansion
  coefficient of the container material
  ($= 3\alpha_{\text{container}}$)
\end{itemize}
\textbf{Physical interpretation:} When the liquid is
heated, both the liquid and the container expand.
The container grows larger, providing more room, so
the liquid's rise in the container understates its
true expansion. What we observe is the excess
expansion of the liquid over and above that of the
container.
\end{lawbox}

\begin{examplebox}[13.4]
\stepcounter{workedexample}
Mercury in a glass flask appears to expand by
$1.72\times10^{-5}\ \text{m}^3$ when heated from
$20\,°\text{C}$ to $120\,°\text{C}$. The original
volume of mercury is $1.00\times10^{-3}\ \text{m}^3$.
($\alpha_{\text{glass}} = 3.2\times10^{-6}\ \text{K}^{-1}$)\\
(a) Calculate the apparent coefficient of expansion
of mercury.\\
(b) Calculate the real coefficient of expansion and
compare it with the value in Table~\ref{tab:expansion}.
\end{examplebox}

\begin{solutionbox}
\textbf{(a)} Apparent coefficient:
\[
  \gamma_{\text{app}}
  = \frac{\Delta V_{\text{apparent}}}{V_0 \Delta\theta}
  = \frac{1.72\times10^{-5}}{1.00\times10^{-3} \times 100}
  = 1.72\times10^{-4}\ \text{K}^{-1}
\]

\textbf{(b)}
$\gamma_{\text{container}} = 3\alpha_{\text{glass}}
= 3 \times 3.2\times10^{-6}
= 9.6\times10^{-6}\ \text{K}^{-1}$

\[
  \gamma_{\text{real}}
  = \gamma_{\text{app}} + \gamma_{\text{container}}
  = 1.72\times10^{-4} + 9.6\times10^{-6}
  = 1.816\times10^{-4}
  \approx 1.82\times10^{-4}\ \text{K}^{-1}
\]
This agrees with the table value of
$182\times10^{-6}\ \text{K}^{-1}$ for mercury.
The true expansion of mercury is slightly larger than
what is observed because the glass container also
expanded (enlarging the space), partially masking
the mercury's expansion. This is precisely the
correction that must be applied when using a
mercury-in-glass thermometer to measure absolute
volumes of expansion.
\end{solutionbox}

% ============================================================
\section{Anomalous Expansion of Water}
% ============================================================

Water behaves unlike all other liquids in one crucial
respect: it does not expand uniformly as it is heated
from $0\,°\text{C}$.

\begin{processbox}[The Anomalous Behaviour of Water]
\begin{itemize}[noitemsep]
  \item From $0\,°\text{C}$ to $4\,°\text{C}$: water
  \textbf{contracts} (density increases)
  \item At $4\,°\text{C}$: water has its
  \textbf{maximum density}
  ($\approx 1000\ \text{kg\,m}^{-3}$)
  \item Above $4\,°\text{C}$: water expands normally
  (density decreases)
  \item Ice is \textbf{less dense} than liquid water
  ($\rho_{\text{ice}} \approx 917\ \text{kg\,m}^{-3}$),
  so ice floats --- unusual for a solid
\end{itemize}
\textbf{Molecular explanation:} In the liquid state,
water molecules are held together by \textit{hydrogen
bonds} --- relatively strong intermolecular
attractions between the partial positive charge on a
hydrogen atom of one molecule and the partial negative
charge on the oxygen atom of a neighbour. Ice has a
rigid, open hexagonal lattice structure maintained by
these hydrogen bonds, with large empty spaces between
molecules; this open structure is less dense than the
disordered liquid. Just above $0\,°\text{C}$, some of
this open ice-like hydrogen-bonded structure persists
in the liquid. As the liquid is warmed from
$0\,°\text{C}$, these residual structures collapse,
packing the molecules more closely together and
causing the water to contract. By $4\,°\text{C}$, the
ice-like clusters have substantially broken down and
the water is at its most compact. Above $4\,°\text{C}$,
normal thermal expansion (increasing molecular
separation due to greater vibrational energy) takes
over, and the density falls as in any ordinary liquid.
\end{processbox}

\begin{diagbox}[Density of Water vs Temperature]
\begin{center}
\begin{tikzpicture}
\begin{axis}[
  width=11cm, height=6cm,
  xlabel={Temperature (\textdegree C)},
  ylabel={Density (kg\,m$^{-3}$)},
  xmin=-5, xmax=25,
  ymin=996.5, ymax=1000.5,
  grid=major, grid style={gray!20},
  xtick={-4,0,4,8,12,16,20,24},
  ytick={997.0,997.5,998.0,998.5,999.0,999.5,1000.0,1000.5},
  clip=true,
]
  % Ice region shading
  \fill[cyan!10] (axis cs:-5,996.5) rectangle (axis cs:0,1000.5);
  % Ice label
  \node[font=\tiny, align=center, cyan!60!black]
    at (axis cs:-2.5,998.4) {Ice\\$\rho < 1000$};
  % Empirical pure-water density fit (kg m^-3) for 0–25°C.
  % It gives a maximum near 4°C and remains inside the plotted axis.
  \addplot[codexblue, very thick, domain=0:25, samples=100]
    {999.842594 + 0.06793952*x - 0.00909529*x^2
      + 0.0001001685*x^3 - 1.120083e-6*x^4 + 6.536332e-9*x^5};
  % Max density marker
  \addplot[only marks, mark=*, codexred, mark size=4pt]
    coordinates {(4, 999.975)};
  \draw[dashed, codexred, thick]
    (axis cs:4,996.5) -- (axis cs:4,999.975);
  \node[codexred, above right, font=\small, align=center]
    at (axis cs:4,999.975) {Max density\\near 4\textdegree C};
  % Annotation arrows (nodes separated from draw so align= works cleanly)
  \draw[->, codexblue, thick]
    (axis cs:1,999.6) -- (axis cs:3.5,999.93);
  \node[above left, font=\tiny, align=center, codexblue]
    at (axis cs:1,999.6) {contracting\\(0 to 4\textdegree C)};
  \draw[->, codexgreen, thick]
    (axis cs:8,999.75) -- (axis cs:14,999.35);
  \node[above right, font=\tiny, align=center, codexgreen]
    at (axis cs:8,999.75) {expanding\\(above 4\textdegree C)};
\end{axis}
\end{tikzpicture}
\end{center}
\end{diagbox}

\begin{realworldbox}[Why Fish Survive Winter]
In a lake during cold weather, the surface water cools
first. As it cools from above $4\,°\text{C}$ toward
$4\,°\text{C}$, it becomes denser and sinks, replaced
by warmer water from below (convection). At
$4\,°\text{C}$, the water is at maximum density and
stops sinking. Further cooling of the surface toward
$0\,°\text{C}$ makes the surface water less dense, so
it stays on top. The surface eventually freezes.
The ice layer acts as an insulator --- ice is a poor
conductor of heat --- and the water beneath, at
$4\,°\text{C}$ or above, remains liquid. Fish and other
aquatic life survive the winter in this liquid water
beneath the ice. Without the anomalous expansion of
water, lakes would freeze solid from the bottom up,
killing all aquatic life.

Beyond ecology, the anomalous expansion has important
geological consequences. Water seeps into cracks in
rocks and, as it freezes, expands by about $9\%$ in
volume. This generates enormous pressures (sufficient
to split granite) and is the primary mechanism of
\textit{frost weathering} or \textit{freeze--thaw
action} --- the process that over geological time
breaks rocks apart, contributes to soil formation,
and shapes mountain landscapes. Water pipes, road
surfaces, and concrete structures are also vulnerable
to freeze--thaw damage for the same reason.
\end{realworldbox}

% ============================================================
\section{Chapter Summary}
% ============================================================

\begin{summarybox}
\textbf{Temperature vs.\ heat:} Temperature is a measure
of average molecular KE (a state property). Heat is
energy in transit between bodies at different
temperatures. They are not the same thing.

\textbf{Zeroth Law:} if A $\leftrightarrow$ C and
B $\leftrightarrow$ C (thermal equilibrium), then
A $\leftrightarrow$ B. Justifies thermometers.

\textbf{Scale conversions:}
$T(\text{K}) = T(°\text{C}) + 273$;
$T(°\text{F}) = \frac{9}{5}T(°\text{C}) + 32$;
$T(\text{K}) = \frac{5}{9}[T(°\text{F}) + 460]$.
Note: write $300\ \text{K}$, \textbf{not}
$300\,°\text{K}$.

\textbf{Calibration:}
$\theta = \frac{X_\theta - X_0}{X_{100}-X_0}
\times 100\,°\text{C}$

\textbf{Linear expansion:}
$\Delta L = \alpha L_0\Delta\theta$
($\alpha$ treated as constant for small $\Delta\theta$)

\textbf{Area expansion:}
$\beta = 2\alpha$

\textbf{Volume expansion:}
$\gamma = 3\alpha$

\textbf{Real vs apparent:}
$\gamma_{\text{real}} = \gamma_{\text{app}}
+ \gamma_{\text{container}}$

\textbf{Bimetallic strip:} two metals with different
$\alpha$ bonded together; curves toward the metal
with the smaller $\alpha$ when heated.
Used in thermostats, circuit breakers, fire alarms.

\textbf{Anomalous water:} maximum density at
$4\,°\text{C}$; ice less dense than water ($\rho_{\text{ice}}
\approx 917\ \text{kg\,m}^{-3}$). Life and geology
depend on this.
\end{summarybox}

% ============================================================
\newpage
\begin{exercisebox}[13]

\subsection*{Section A: Multiple Choice}

\textbf{A1.} The temperature $-40\,°\text{C}$ in
Kelvin is:

\mcoption{A}{$233\ \text{K}$}
\mcoption{B}{$-313\ \text{K}$}
\mcoption{C}{$313\ \text{K}$}
\mcoption{D}{$40\ \text{K}$}
\ansref{Ch13-A1}

\vspace{0.4cm}

\textbf{A2.} A mercury thermometer has length
$2.0\ \text{cm}$ at $0\,°\text{C}$ and
$12.0\ \text{cm}$ at $100\,°\text{C}$. At
$25\,°\text{C}$, the length is:

\mcoption{A}{$3.0\ \text{cm}$}
\mcoption{B}{$4.5\ \text{cm}$}
\mcoption{C}{$4.8\ \text{cm}$}
\mcoption{D}{$5.0\ \text{cm}$}
\ansref{Ch13-A2}

\vspace{0.4cm}

\textbf{A3.} A steel rod has length $2.0\ \text{m}$
at $20\,°\text{C}$. At $70\,°\text{C}$, its length is
($\alpha_{\text{steel}} = 12\times10^{-6}\ \text{K}^{-1}$):

\mcoption{A}{$2.0012\ \text{m}$}
\mcoption{B}{$2.00012\ \text{m}$}
\mcoption{C}{$2.0120\ \text{m}$}
\mcoption{D}{$2.0006\ \text{m}$}
\ansref{Ch13-A3}

\vspace{0.4cm}

\textbf{A4.} The coefficient of volume expansion
$\gamma$ is related to the linear expansion
coefficient $\alpha$ by:

\mcoption{A}{$\gamma = \alpha$}
\mcoption{B}{$\gamma = 2\alpha$}
\mcoption{C}{$\gamma = 3\alpha$}
\mcoption{D}{$\gamma = \alpha^3$}
\ansref{Ch13-A4}

\vspace{0.4cm}

\textbf{A5.} Water has its maximum density at:

\mcoption{A}{$0\,°\text{C}$}
\mcoption{B}{$4\,°\text{C}$}
\mcoption{C}{$10\,°\text{C}$}
\mcoption{D}{$100\,°\text{C}$}
\ansref{Ch13-A5}

\vspace{0.4cm}

\textbf{A6.} A bimetallic strip curves when heated
because:

\mcoption{A}{Both metals expand by the same amount}
\mcoption{B}{The two metals have different
coefficients of linear expansion}
\mcoption{C}{One metal melts before the other}
\mcoption{D}{The strip loses mass on heating}
\ansref{Ch13-A6}

\vspace{0.4cm}

\textbf{A7.} The apparent expansion of a liquid is
\textit{less than} the real expansion because:

\mcoption{A}{Some liquid evaporates}
\mcoption{B}{The container also expands}
\mcoption{C}{Gravity compresses the liquid}
\mcoption{D}{The liquid contracts when heated}
\ansref{Ch13-A7}

\vspace{0.4cm}

\textbf{A8.} Which temperature scale has its zero
point at absolute zero and is the SI base unit of
thermodynamic temperature?

\mcoption{A}{Celsius}
\mcoption{B}{Fahrenheit}
\mcoption{C}{Kelvin}
\mcoption{D}{Newton scale}
\ansref{Ch13-A8}

\vspace{0.4cm}

\textbf{A9.} A circular hole in a metal plate is
heated. The hole:

\mcoption{A}{Gets smaller (metal expands inward)}
\mcoption{B}{Gets larger (the metal around it
expands outward)}
\mcoption{C}{Stays the same size}
\mcoption{D}{Depends on whether the hole is in
the centre or edge of the plate}
\ansref{Ch13-A9}

\vspace{0.4cm}

\textbf{A10.} The Zeroth Law of Thermodynamics
justifies:

\mcoption{A}{The impossibility of perpetual motion}
\mcoption{B}{The use of thermometers to measure
temperature}
\mcoption{C}{The fact that heat flows from hot to cold}
\mcoption{D}{The expansion of gases on heating}
\ansref{Ch13-A10}

\vspace{0.6cm}

\subsection*{Section B: Theory and Calculations}

\textbf{B1.} A concrete bridge span is $120\ \text{m}$
long at $15\,°\text{C}$. The bridge must withstand
temperatures ranging from $-5\,°\text{C}$ (cold night)
to $55\,°\text{C}$ (hot day).
($\alpha_{\text{concrete}} = 12\times10^{-6}\ \text{K}^{-1}$)\\[4pt]
(a) Calculate the total change in length over this
temperature range.\\
(b) Expansion joints are placed at each end of the
span. How wide must each joint be to accommodate the
full expansion?\\
(c) Modern bridges use prestressed concrete with
steel reinforcement. If a steel rod of length
$120\ \text{m}$
($\alpha_{\text{steel}} = 12\times10^{-6}\ \text{K}^{-1}$)
is embedded in the concrete and both heat up together,
explain why there is no cracking due to differential
expansion.\\
(d) If a glass bridge beam ($\alpha_{\text{glass}}
= 3.2\times10^{-6}\ \text{K}^{-1}$) were embedded
instead, what problem would arise? Calculate the
differential expansion.
\hfill\ansref{Ch13-B1}

\vspace{0.4cm}

\textbf{B2.} A resistance thermometer has resistance
$50.0\ \Omega$ at $0\,°\text{C}$ and $68.4\ \Omega$
at $100\,°\text{C}$.
A thermocouple thermometer simultaneously reads
$45.0\,°\text{C}$ when the resistance thermometer
reads $58.3\ \Omega$.\\[4pt]
(a) Calculate the temperature indicated by the
resistance thermometer.\\
(b) Explain why the two readings differ.\\
(c) The resistance at a very high temperature is
$130.0\ \Omega$. Calculate the temperature on the
platinum resistance scale.\\
(d) Explain why gas thermometers are used as
primary standards rather than liquid-in-glass or
resistance thermometers.
\hfill\ansref{Ch13-B2}

\vspace{0.4cm}

\textbf{B3.} A glass bottle of volume
$500\ \text{cm}^3$ is filled completely with
glycerol at $20\,°\text{C}$.
($\gamma_{\text{glycerol}} = 4.85\times10^{-4}\ \text{K}^{-1}$,
$\alpha_{\text{glass}} = 3.2\times10^{-6}\ \text{K}^{-1}$)\\[4pt]
(a) Calculate the volume expansion coefficient of
the glass bottle.\\
(b) When heated to $80\,°\text{C}$, calculate the
volume of glycerol that overflows.\\
(c) Calculate the apparent coefficient of expansion
of glycerol in glass.\\
(d) Why does liquid overflow from a completely full
container when heated, even though both the liquid
and container expand?
\hfill\ansref{Ch13-B3}

\vspace{0.4cm}

\textbf{B4.} (a) Draw a careful sketch of the
density-temperature graph for water from
$-2\,°\text{C}$ to $20\,°\text{C}$, labelling the
key features.\\[4pt]
(b) Explain at the molecular level why ice is less
dense than liquid water.\\
(c) A lake is $10\ \text{m}$ deep and cools uniformly.
Describe, step by step, the convection process that
occurs as the lake cools from $10\,°\text{C}$ toward
$0\,°\text{C}$, and explain why the bottom of the
lake remains at $4\,°\text{C}$ even when the surface
is frozen.\\
(d) Suggest two consequences of the anomalous
expansion of water for life on Earth, beyond the
fish survival example.
\hfill\ansref{Ch13-B4}

\vspace{0.4cm}

\textbf{B5.} An aluminium sphere of diameter
$5.0\ \text{cm}$ has a hole of diameter
$2.0\ \text{cm}$ drilled through its centre. The
sphere is heated from $20\,°\text{C}$ to $200\,°\text{C}$.
($\alpha_{\text{Al}} = 23\times10^{-6}\ \text{K}^{-1}$)\\[4pt]
(a) Calculate the new outer diameter of the sphere.\\
(b) Calculate the new diameter of the hole.\\
(c) Calculate the new volume of the sphere
(including the hole volume).\\
(d) A steel rod of diameter $2.002\ \text{cm}$ at
$20\,°\text{C}$ is pushed into the hole.
($\alpha_{\text{steel}} = 12\times10^{-6}\ \text{K}^{-1}$)\\
At what temperature will the rod just fit through
the hole? (Hint: set the hole diameter equal to the
rod diameter and solve for $T$.)
\hfill\ansref{Ch13-B5}

\end{exercisebox}
